% BOOK-FILE-ID: f-4663a7fcd21b
% subsection: 対数型 \, $a_{n+1}=pa_n^q$

\subsection{\texorpdfstring{対数型  $a_{n+1}=pa_n^q$}{対数型}}\label{節：対数型}
  今までは$a_n$に対して$f(n)$を足したり掛けたりしてきたが,\,$a_n$自体を掛けたり足したりしたらどうなるだろう.

  まずはシンプルに
  \[
    a_{n+1}=pa_n^q
  \]
  という漸化式を考えてみることにする.
  さて,いかにして既知の形に帰着させようか.
  $n$倍であったり,\,$n$を足すことに関してはここまでの内容を押さえた読者ならばできることだろう.
  しかし$q$乗をするという操作については今までしたことがない.
  この$q$乗という操作を足したり掛けたりすることにできれば解くことができそうだ.
  ここで1つの知識を思い出そう.
  対数は積を和に,累乗を定数倍に変えるものだった.
  だから,今回のように累乗を掛け算に変換するのにはうってつけの道具であろう.\footnote{ \sectionref*{節：指数と対数}で対数について扱っている.対数に自信がない人は参照すると良い.}

  ここで1つ断りをいれておこう.
  この節の$q$は指数を表し,特性方程式型で定数項を表した$q$とは別である.
  また$p>0$とし,対数の真数条件から,この節で扱う$a_n$はすべて0より大きい実数であるとする.
  以後断りなしに対数をとるが,答案等に書く際は常に真数条件を確かめる\footnote{$a_n>0$ならば$p a_n^q>0$であるから,\bookterm{induction}{数学的帰納法}によりすべての項が正となり,\bookterm{logarithm}{対数}を取ることができる.}こと.
  

  まず$p\ne1$として,両辺の底$p$の対数を取ってみると,
  \begin{align*}
      \log_pa_{n+1}&=\log_pp+q\log_pa_n\\
      &=q\log_pa_n+1.
  \end{align*}
  ここで底を$p$とする対数をとったが,底は正で1以外なら何を選んでもよい.
  $p=1$のときは底を$p$にできないので,例えば$b_n=\log_2a_n$と置くと$b_{n+1}=qb_n$となる.
  以下では引き続き$p\ne1$の場合を考える.
  $b_n=\log_pa_n$とおくと
  \[
    b_{n+1}=qb_n+1
  \]
  を得る.
  対数を取った先では,\,$q$乗という見慣れない操作が,\,$q$倍という既知の操作に変わったのだ.
  よってこの漸化式は特性方程式型になったので,これを解くと
  \[
    b_n=\left(b_1-\frac{1}{1-q}\right)q^{n-1}+\frac{1}{1-q}
  \]
  となる.ただし$b_1=\log_pa_1$であり,この式は$q\ne0,1$の場合である.\footnote{$q=1$なら等差型に帰着する.\,$q=0$なら元の式から$n\ge2$で$a_n=p$と分かる.}
  したがって
  \[
    b_n=\log_p a_n  \iff a_n=p^{b_n}
  \]
  であることから,一般項を得る.

  また,階比型漸化式を対数を使って解くことで,より楽に解ける場合がある.
  それは,積を和に変えるという対数の便利な性質に由来するものである.
  
  \noindent\bookblocklink{example-logarithmic-1}{problem}{answer}{【例題】}\par\nobreak
  \begin{enumerate}[label=(\arabic*),parsep=3pt,format=\bookitemlink{example-logarithmic-1}{problem}{answer}]
    \item $a_1=2,\quad a_{n+1}=2a_n^2$
    \item $a_1=\frac{1}{2},\quad a_{n+1}=\frac{1}{a_n}$
    \item $a_1=1,\quad a_{n+1}=2^na_n$
    \item $a_1=5,\quad a_{n+1}=a_n^2-2a_n+2$
    \item $a_1=2,\quad a_{n+1}=\frac{a_n^2}{2^n}$
    \item $a_1=1,\quad a_{n+1}=(n+1)a_n$
  \end{enumerate}
  
  \noindent\bookblocklink{example-logarithmic-1}{hint}{problem}{【方針】}\par\nobreak
  \begin{enumerate}[label=(\arabic*),parsep=3pt,format=\bookitemlink{example-logarithmic-1}{hint}{problem}]
    \item $b_n=\log_2a_n$で$b_{n+1}=2b_n+1$に変え,定数項を消します.
    \item 逆数を\bookterm{logarithm}{対数}で符号反転に変えます.
    \item $f(n)=2^n$が指数関数なので,\bookterm{logarithm}{対数}を取ると和の形に帰着できます.
    \item \bookterm{square}{平方完成}で$a_{n+1}-1=(a_n-1)^2$を作ります.\bookterm{logarithm}{対数}を取る量は$a_n$ではなく$a_n-1$であり,その正値を先に確認します.
    \item $b_n=\log_2a_n$から一次式$n+1$を引くと等比型になります.
    \item 積を直接まとめる方法と,\bookterm{logarithm}{対数}で和に変える方法を比べましょう.最後には同じ積$n!$を計算します.
  \end{enumerate}
  
  \noindent\bookblocklink{example-logarithmic-1}{answer}{problem}{【解答】}\par\nobreak
  \begin{enumerate}[label=(\arabic*),parsep=3pt,format=\bookitemlink{example-logarithmic-1}{answer}{problem}]
    \item $a_n=2^{2^{n}-1}$
    \item $a_n=2^{(-1)^{n}}$
    \item $a_n=2^{\frac{n(n-1)}{2}}$
    \item $a_n=2^{2^n}+1$
    \item $a_n=2^{\,n+1-2^{n-1}}$
    \item $a_n=n!$
  \end{enumerate}
  
  \noindent\bookblocklink{example-logarithmic-1}{solution}{problem}{【解法】}\par\nobreak
  
  (1),(2),(3),(5),(6)は初項が正で,正の項から次も正になる.\bookterm{induction}{数学的帰納法}より全項が正である.
  (4)は$a_1-1=4>0$と$a_{n+1}-1=(a_n-1)^2$から,\,$a_n-1>0$を保つ.この量に\bookterm{logarithm}{対数}を取る.
  \begin{enumerate}[label=(\arabic*),parsep=3pt,format=\bookitemlink{example-logarithmic-1}{solution}{problem}]
    \item $b_n=\log_2a_n$と置くと,
      \[
        b_{n+1}=2b_n+1,\qquad b_1=1.
      \]
      $b_n+1$は初項$2$,公比$2$の等比型なので,
      \[
        b_n+1=2^n,\qquad a_n=2^{b_n}=2^{2^n-1}.
      \]
    \item $b_n=\log_2a_n$と置くと,逆数は
      \[
        b_{n+1}=-b_n,\qquad b_1=\log_2\frac12=-1
      \]
      に変わる.従って$b_n=(-1)^n$であり,\,$a_n=2^{(-1)^n}$となる.
    \item 自然\bookterm{logarithm}{対数}を取れば$\log a_{n+1}=\log a_n+n\log2$である.差を足すと,
      \[
        \log a_n=\sum_{k=1}^{n-1}k\log2=\frac{n(n-1)}2\log2.
      \]
      従って$a_n=2^{\dfrac{n(n-1)}{2}}$となる.
    \item \bookterm{square}{平方完成}で$a_{n+1}-1=(a_n-1)^2$となる.\bookterm{logarithm}{対数}を取ると,
      \[
        \log_2(a_{n+1}-1)=2\log_2(a_n-1).
      \]
      初項は$\log_2(a_1-1)=2$だから,\,$\log_2(a_n-1)=2^n$であり,\,$a_n=2^{2^n}+1$となる.
    \item $b_n=\log_2a_n$と置くと,\,$b_{n+1}=2b_n-n,\ b_1=1$である.
      一次式$n+1$を引けば,
      \[
        b_{n+1}-(n+2)=2\{b_n-(n+1)\}.
      \]
      $c_n=b_n-(n+1)$は$c_{n+1}=2c_n,\ c_1=-1$を満たすので,
      \[
        c_n=-2^{n-1},\quad b_n=n+1-2^{n-1}.
      \]
      戻すと$a_n=2^{\,n+1-2^{n-1}}$となる.
    \item 自然\bookterm{logarithm}{対数}を取れば$\log a_{n+1}=\log a_n+\log(n+1)$である.
      \[
        \log a_n=\log1+\log2+\cdots+\log n=\log n!.
      \]
      従って$a_n=n!$となる.直接積を取っても同じ結果が得られる.
  \end{enumerate}
  \paragraphbf{補足}
  
  (3),(6)の$\log$は底が$e$の自然\bookterm{logarithm}{対数}である.
  変形後の真数が$0$や負なら,実数の\bookterm{logarithm}{対数}は取れない.例えば(4)で初項が$1$なら定数\bookterm{sequence}{数列}$a_n=1$となり,\,$\log(a_n-1)$を取らずに扱う.
  
  \noindent\bookblocklink{check-log-domain}{problem}{answer}{【確認問題】}\par\nobreak
  \begin{enumerate}[format=\bookitemlink{check-log-domain}{problem}{answer}]
    \item $a_{n+1}=a_n^2$で,初項が$2,0,-2$の場合を比較せよ.
      最初から実数の\bookterm{logarithm}{対数}を取れるのはどの場合か.取れない場合にも\bookterm{sequence}{数列}は定義されるか.
  \end{enumerate}
  \noindent\bookblocklink{check-log-domain}{answer}{problem}{【確認の解答】}\par\nobreak
  \begin{enumerate}[format=\bookitemlink{check-log-domain}{answer}{problem}]
    \item 初項$2$ならすべての項が正である.
      初項$0$なら全項が0の数列,初項$-2$なら第2項以降が正である.
      いずれも元の\bookterm{recurrence}{漸化式}は定義される.\bookterm{logarithm}{対数}を取る条件と,元の\bookterm{sequence}{数列}が存在する条件は区別する.
  \end{enumerate}
